\begin{thebibliography}{8}

\bibitem{ref1}
R. A. Yeh, A. G. Schwing, J. Huang, and K. Murphy, ``Diverse generation for multi-agent sports games,'' in Proc. IEEE/CVF Conf. Comput. Vis. Pattern Recognit. (CVPR), 2019.

\bibitem{ref2}
W. Z. Teoh, ``Coherent multi-agent trajectory forecasting in team sports with CausalTraj,'' arXiv preprint arXiv:2511.18248, 2025.

\bibitem{ref3}
FOOTPASS Research Team, ``FOOTPASS: A multi-modal multi-agent tactical context dataset for play-by-play action spotting in soccer broadcast videos,'' arXiv preprint arXiv:2511.16183, 2025.

\bibitem{ref4}
K. Kurach et al., ``Google research football: A novel reinforcement learning environment,'' in Proc. AAAI Conf. Artif. Intell., 2020.

\bibitem{ref5}
B. Brandao, T. W. De Lima, A. Soares, and L. Melo, ``Multi-agent reinforcement learning for tactical decision-making in soccer,'' ResearchGate, 2022.

\bibitem{ref6}
D. Gupta, ``Autonomous agentic AI frameworks for supply chain resilience and decision-making,'' in Proc. 5th Int. Conf. Innov. Practices Technol. Manage. (ICIPTM), 2026.

\bibitem{ref7}
T. Hazenberg et al., ``Multi-agent reinforcement learning for dynamic pricing in supply chains: Benchmarking strategic agent behaviours under realistically simulated market conditions,'' arXiv preprint arXiv:2507.02698, 2025.

\bibitem{ref8}
M. Bhatnagar, ``FanCric: Multi-agentic framework for crafting fantasy 11 cricket teams,'' arXiv preprint arXiv:2410.01307, 2024.

\end{thebibliography}